\scanpage{6}
une \deleted{autre} \inserted{expression de l’}accouplement $*$ en termes
de $(X,\underline a)_v$. \deleted{On voit \unsure{s|q}}

\[
 \text{\uline{Cor}}\quad\unsure{v|\nu}(X_{\underline a},\underline b)=v(X_{\underline b}^{-},\underline a)\quad(=v(X_{b^{-}},\underline a^{-}))
\]

Dans le cas global [où on a une
famille de \uline{valuations} \uline{discrètes} satisfaisant la
formule du produit], on obtient un
accouplement\deleted{\unsure{aux}} $D_a^{\unsure{*}}(A_K)\times\deleted{\unsure{D}}Z_0^*(A_K)\longrightarrow\mathbb R$
par \deleted{produit} produit des symboles locaux

\[
 (X,\underline a)=\sum_v(X,\underline a)_v
\]

qui ne dépend plus que des classes \deleted{de}
mod $D_\ell(A_K)$ resp. mod $Z_\ell^*(A_K)$, et définit
donc un accouplement

\[
 A^*(K)\times A^{\deleted{\unsure{*}}}(K)\longrightarrow\mathbb R .
\]

On obtient donc : cet accouplement est
à valeurs dans $\frac1m\mathbb Z\subset\mathbb Q$, où
$m$ est un entier $>0$ convenable

La forme en question est la \uline{forme de
Néron-Tate}. \deleted{Elle est \unsure{multipliée|multiplicative}}
% archive sheet number: 5
